\chapter{距離空間と位相}
\label{ch:tb-topology}

位相では定義そのものが答案の初手になるため、ここだけは開・閉・コンパクト・連結を正確に言えるようにする。
証明全文を暗記する必要はなく、「開被覆を逆像へ戻す」「二点の間で切る」などの操作を再現する。

\section{開集合・閉集合・閉包}

\begin{definition}{距離空間と開球}{tb-metric-space}
集合 $X$ 上の距離 $d$ は
\[
d(x,y)\ge0,\quad d(x,y)=0\Longleftrightarrow x=y,\quad
d(x,y)=d(y,x),\quad d(x,z)\le d(x,y)+d(y,z)
\]
を満たす。$B(x;r)=\{y\in X\mid d(x,y)<r\}$ を開球という。
\end{definition}

\begin{definition}{開集合・閉集合・閉包}{tb-open-closed-closure}
$U\subset X$ が開集合であるとは、各 $x\in U$ に対し
$B(x;r)\subset U$ となる $r>0$ が存在することをいう。
$F\subset X$ が閉集合であるとは $X\setminus F$ が開集合であることをいう。
$A$ を含む最小の閉集合を閉包 $\overline A$ といい、
$x\in\overline A$ を $A$ の触点という。
\end{definition}

距離空間では
\[
x\in\overline A
\Longleftrightarrow \forall r>0,\ B(x;r)\cap A\ne\varnothing
\Longleftrightarrow A\text{ の点列 }a_n\to x\text{ が存在する}.
\]
$\overline A=X$ のとき $A$ は $X$ で稠密という。

閉包を求める問題では、まず集合内の典型的な点列の極限を探す。
閉集合を示すときは、補集合が開であることと、収束列の極限を含むことの短い方を選ぶ。

\begin{proposition}{閉包の計算法則}{tb-closure-laws}
\[
A\subset\overline A,\quad
\overline{\overline A}=\overline A,\quad
A\subset B\Rightarrow\overline A\subset\overline B,\quad
\overline{A\cup B}=\overline A\cup\overline B.
\]
\end{proposition}

最後の式は、右辺が $A\cup B$ を含む閉集合なので一方の包含を得る。
逆包含は $A,B\subset A\cup B$ の単調性から得る。

\begin{definition}{相対位相}{tb-subspace-topology}
$Y\subset X$ に対し、$Y$ の開集合を $Y\cap U$（$U$ は $X$ の開集合）と定める。
閉集合も $Y\cap F$（$F$ は $X$ の閉集合）の形である。
したがって $[1,2]\cap\Q$ は $\Q$ の中では閉集合である。
\end{definition}

\begin{definition}{位相空間と Hausdorff 性}{tb-topological-space}
集合 $X$ の部分集合族 $\mathcal T$ が
\[
\varnothing,X\in\mathcal T,\qquad
\mathcal T\text{ の任意個の元の合併が }\mathcal T\text{ に属する},
\]
\[
\mathcal T\text{ の有限個の元の共通部分が }\mathcal T\text{ に属する}
\]
を満たすとき $(X,\mathcal T)$ を位相空間という。$\mathcal T$ の元を開集合、補集合を閉集合と呼ぶ。
異なる二点が互いに交わらない開近傍で分離できるとき Hausdorff 空間という。
\end{definition}

\begin{definition}{点列の収束}{tb-topological-sequence}
位相空間 $X$ の点列 $a_n$ が $a\in X$ に収束するとは、
$a$ の任意の開近傍 $U$ に対し、ある $N$ が存在して
\[
n\ge N\Longrightarrow a_n\in U
\]
となることをいう。
\end{definition}

\begin{proposition}{Hausdorff 空間では極限が一意}{tb-hausdorff-limit}
Hausdorff 空間の点列が極限を持つなら、その極限は一つだけである。
\end{proposition}

\begin{proof}
$a_n\to a,b$ かつ $a\ne b$ と仮定する。Hausdorff 性により
$a\in U$, $b\in V$, $U\cap V=\varnothing$ となる開近傍を取れる。
十分大きい $n$ では $a_n\in U$ かつ $a_n\in V$ となり矛盾する。
\end{proof}

\begin{proposition}{補有限位相}{tb-cofinite-topology}
無限集合 $X$ に
\[
\mathcal T=\{\varnothing\}\cup
\{U\subset X\mid X\setminus U\text{ は有限集合}\}
\]
を入れると位相空間になり、コンパクトかつ連結である。
\end{proposition}

任意個の合併の補集合は有限集合の一つに含まれ、有限個の共通部分の補集合は
有限集合の有限合併なので、位相の公理を満たす。
開被覆から非空な一つ $U_0$ を選べば $X\setminus U_0$ は有限であり、
残る各点を覆う開集合を一つずつ足せば有限部分被覆になる。
また二つの非空開集合の共通部分は補有限、従って空でない。
よって二つの非空開集合による分離は存在しない。

\section{連続写像}

\begin{definition}{連続性}{tb-metric-continuity}
距離空間の写像 $f:X\to Y$ が $x\in X$ で連続であるとは
\[
\forall\varepsilon>0\ \exists\delta>0\ \forall u\in X,\qquad
d_X(u,x)<\delta\Rightarrow d_Y(f(u),f(x))<\varepsilon
\]
をいう。距離空間では $x_n\to x\Rightarrow f(x_n)\to f(x)$ と同値である。
\end{definition}

\begin{theorem}{開集合・閉集合による連続性}{tb-continuity-preimage}
$f:X\to Y$ について次は同値である。
\[
f\text{ は連続},\quad
Y\text{ の全ての開集合 }U\text{ で }f^{-1}(U)\text{ は開},\quad
Y\text{ の全ての閉集合 }F\text{ で }f^{-1}(F)\text{ は閉}.
\]
\end{theorem}

補集合と逆像が交換する
$f^{-1}(Y\setminus F)=X\setminus f^{-1}(F)$ ため、開集合版と閉集合版は同値である。

\begin{proposition}{閉包と連続写像}{tb-continuity-closure}
連続写像 $f:X\to Y$ と $A\subset X$ に対し
\[
f(\overline A)\subset\overline{f(A)}.
\]
\end{proposition}

\begin{proof}
$\overline{f(A)}$ は閉集合なので、その逆像は閉集合である。
しかも $A\subset f^{-1}(\overline{f(A)})$ だから
$\overline A\subset f^{-1}(\overline{f(A)})$、従って主張を得る。
\end{proof}

$\abs{f(x)-f(y)}\le Ld(x,y)$ を満たす写像は Lipschitz 連続である。
この評価が得られたら $\delta=\varepsilon/L$ と取ればよい。
具体的な距離評価があれば Lipschitz、点列が与えられれば点列、
位相だけの問題なら開集合の逆像を使うのが最短である。

\begin{theorem}{コンパクト集合上では一様連続}{tb-heine-cantor}
コンパクト距離空間 $K$ 上の連続写像は一様連続である。
\end{theorem}

この定理は「点ごとの連続性」を「全ての点で共通の $\delta$」へ変える。
例えば $I=[0,1]$ とし、連続関数 $g:I\times I\to\R$ が
$c=\inf_{x\in I}g(x,0)>0$ を満たすとする。一様連続性から、十分小さい $b>0$ では全ての $x$ で
\[
|g(x,b)-g(x,0)|<c/2,\qquad g(x,b)>c/2.
\]
従って $\inf_xg(x,b)=0$ にはなれない。一方、$\R\times I$ はコンパクトでないため
\[
f(x,b)=(bx-1)^2
\]
なら $\inf_xf(x,0)=1$ だが、$b>0$ では $x=1/b$ で最小値0を取る。

\begin{definition}{商位相}{tb-quotient-topology}
同値関係 $\sim$ による商集合 $Y=X/{\sim}$ と標準全射 $\pi:X\to Y$ に対し
\[
U\subset Y\text{ が開}\Longleftrightarrow\pi^{-1}(U)\text{ が }X\text{ で開}
\]
と定める。この位相では $\pi$ は定義から連続である。
\end{definition}

$S\subset X$ がコンパクトで $\pi(S)=Y$ まで示せたなら、
$\pi|_S:S\to Y$ は連続全射なので、$Y$ はコンパクト集合の連続像としてコンパクトである。

\section{コンパクト性}

\begin{definition}{コンパクト}{tb-compact}
$K$ の任意の開被覆 $K\subset\bigcup_{\alpha\in I}U_\alpha$ から
有限個 $U_{\alpha_1},\ldots,U_{\alpha_m}$ でなお $K$ を覆えるとき、$K$ はコンパクトという。
\end{definition}

距離空間では「任意の点列が収束部分列を持ち、その極限が $K$ に属する」と同値である。
$\R^n$ では
\[
K\text{ がコンパクト}\Longleftrightarrow K\text{ が有界かつ閉}
\]
である。

\begin{proposition}{点列コンパクトから有限部分被覆を作る}{tb-sequential-compact-proof}
距離空間 $X$ が点列コンパクトなら、任意の開被覆 $\{W_\lambda\}$ に対して
ある $\delta>0$ が存在し、各 $x\in X$ で
\[
B(x;\delta)\subset W_\lambda
\]
となる $\lambda$ を選べる。さらに $X$ はコンパクトである。
\end{proposition}

最初の主張が偽なら、各 $n$ で「$B(x_n;1/n)$ がどの $W_\lambda$ にも入らない」
点 $x_n$ を取れる。収束部分列 $x_{n(k)}\to x$ を取り、$x\in W_\lambda$ を一つ選ぶ。
$W_\lambda$ は開なので $B(x;\varepsilon)\subset W_\lambda$ となる $\varepsilon>0$ がある。
十分大きい $k$ では
\[
d(x_{n(k)},x)<\varepsilon/2,\qquad 1/n(k)<\varepsilon/2
\]
だから $B(x_{n(k)};1/n(k))\subset W_\lambda$ となり矛盾する。

次に $X$ は任意の $r>0$ に対して有限個の半径 $r$ の球で覆える。
そうでなければ、既に選んだ点から距離 $r$ 以上の点を順に選び、
収束部分列を持たない点列を作れてしまう。
上で得た $\delta$ に対して有限個の球 $B(x_i;\delta/2)$ で $X$ を覆い、
各 $i$ で $B(x_i;\delta)\subset W_{\lambda_i}$ を選べば、
$W_{\lambda_1},\ldots,W_{\lambda_m}$ が有限部分被覆になる。

\begin{theorem}{コンパクト性の基本三定理}{tb-compact-toolkit}
\begin{enumerate}
  \item コンパクト空間の閉部分集合はコンパクトである。
  \item コンパクト集合の連続像はコンパクトである。
  \item Hausdorff 空間のコンパクト部分集合は閉集合である。
\end{enumerate}
\end{theorem}

問題が $\R^n$ の部分集合なら最初に「有界かつ閉」を試す。
一般の位相空間なら勝手に点列へ置き換えず、開被覆の定義から始める。

証明の核は次の通りである。
\begin{enumerate}
  \item 閉部分集合 $F\subset K$ の開被覆に $K\setminus F$ を一つ加えて $K$ を覆う。
  \item $f(K)$ の開被覆を逆像で $K$ の開被覆へ戻し、有限個を選んで再び像へ送る。
  \item $x\notin K$ と各 $y\in K$ を Hausdorff 性で分離し、$K$ 側の有限部分被覆を選ぶ。
\end{enumerate}

\begin{corollary}{コンパクトから Hausdorff への単射}{tb-compact-bijection}
$K$ がコンパクト、$Y$ が Hausdorff で、$f:K\to Y$ が連続単射なら
$f:K\to f(K)$ は同相写像である。特に $f$ が全単射なら $K$ と $Y$ は同相である。
\end{corollary}

実際、閉集合 $F\subset K$ はコンパクト、その像 $f(F)$ はコンパクト、従って $Y$ で閉である。
よって $f$ は閉写像であり、$(f^{-1})^{-1}(F)=f(F)$ から逆写像が連続となる。

定義域のコンパクト性を外すと結論は壊れる。例えば
\[
f:[0,2\pi)\to S^1,\qquad f(t)=(\cos t,\sin t)
\]
は連続全単射だが、$(1,0)$ の近くで逆写像の値が $0$ 側と $2\pi$ 側に分かれるため同相写像でない。
Hausdorff 性も必要である。二点集合上の離散位相から密着位相への恒等写像は
コンパクト空間からの連続全単射だが、逆写像は連続でない。

Hausdorff 空間では一点集合が閉である。したがって、コンパクト空間 $S$ から
Hausdorff 空間 $T$ への連続写像 $f$ と $p\in T$ に対し、
$f^{-1}(\{p\})$ は $S$ の閉部分集合、従ってコンパクトである。

\section{連結性と道連結性}

\begin{definition}{連結・道連結}{tb-connected}
$X$ が二つの互いに素な空でない開集合の合併に分かれないとき連結という。
任意の $x,y\in X$ に対し連続写像 $\gamma:[0,1]\to X$ で
$\gamma(0)=x$, $\gamma(1)=y$ となるものが存在するとき道連結という。
道連結なら連結である。
\end{definition}

\begin{theorem}{連続像は連結性を保つ}{tb-connected-image}
$X$ が連結で $f:X\to Y$ が連続なら $f(X)$ は連結である。
$f$ が全射なら特に $Y$ が連結である。連続像は道連結性も保つ。
\end{theorem}

もし $f(X)=U\cup V$ が分離なら、逆像 $f^{-1}(U),f^{-1}(V)$ が $X$ の分離となるため矛盾する。
実数の部分集合が連結であることは区間であることと同値である。
連結を示すなら道を作るか連続像にする。非連結を示すなら、二つの非空な相対開集合を具体的に書く。

\begin{proposition}{二重周期関数}{tb-periodic-function}
連続関数 $f:\R^2\to\R$ が
$f(x+1,y)=f(x,y+1)=f(x,y)$ を満たすなら、最大値 $b$ と最小値 $a$ を持ち、
$f(\R^2)=[a,b]$ である。
\end{proposition}

\begin{proof}
整数部分を引けば任意の点を $[0,1]^2$ へ移せるので
$f(\R^2)=f([0,1]^2)$ である。右辺はコンパクト集合の連続像だからコンパクトで、最大・最小を取る。
それらを取る点を $p,q$ とし $\sigma(t)=(1-t)p+tq$ とおけば、
$f\circ\sigma$ は連続で端点値が $a,b$ だから中間値の定理により $[a,b]$ へ全射である。
\end{proof}

\section{有理数と数列集合の標準例}

\begin{proposition}{有理数の部分空間}{tb-rational-examples}
通常の距離を入れた $\Q$ で $A=[1,2]\cap\Q$ とすると、$A$ は閉だがコンパクトでも連結でもない。
また $\Q$ と $\Q^2$ は全不連結であり、$\R^2\setminus\Q^2$ は道連結である。
\end{proposition}

\begin{proof}
$A=\Q\cap[1,2]$ なので相対位相で閉である。
$\sqrt2$ に収束する $A$ 内の有理数列は $A$ 内に収束部分列を持たないのでコンパクトでない。
二つの異なる有理点の間に無理数を選び、その点で左右に切れば連結でない。
同じ切断により $\Q$ の二点以上を含む部分集合は全て非連結である。

$\Q^2$ の異なる二点も、異なる座標の間を無理数で切れるため、二点以上の連結部分集合を持たない。
$p,q\in\R^2\setminus\Q^2$ に対し、各 $a\in\Q^2$ を通る直線 $pa,qa$ を全て避ける点 $z$ を選ぶ。
避ける直線は可算本しかなく平面を覆わない。このとき折れ線 $p\to z\to q$ は $\Q^2$ を通らない。
\end{proof}

\begin{proposition}{$\{1/n\}$ とその極限点}{tb-sequence-set}
\[
A=\{1/n\mid n\in\N\},\qquad B=A\cup\{0\}
\]
とすると、$A$ は閉でもコンパクトでもなく、$B$ はコンパクトである。
どちらも二点以上を含む連結部分集合を持たない。
\end{proposition}

$0\in\overline A\setminus A$ なので $A$ は閉でない。
開被覆 $U_N=(1/(N+1),2)$（$N\in\N$）は $A$ を覆うが有限部分被覆を持たない。
$B$ の任意の開被覆では、$0$ を含む一つの開集合が $1/n$ の十分長い尾部を覆い、
残る有限個を有限個の開集合で覆える。

\section{集合までの距離}

\begin{definition}{集合までの距離}{tb-distance-to-set}
空でない $A\subset X$ に対し
\[
d(x,A)=\inf_{a\in A}d(x,a)
\]
と定める。
\end{definition}

\begin{theorem}{集合までの距離の三性質}{tb-distance-properties}
\begin{enumerate}
  \item $\abs{d(x,A)-d(y,A)}\le d(x,y)$。従って $x\mapsto d(x,A)$ は連続である。
  \item $d(x,A)=0\Longleftrightarrow x\in\overline A$。
  \item $A$ がコンパクトなら、各 $x$ に対して $d(x,A)=d(x,a_0)$ となる $a_0\in A$ が存在する。
\end{enumerate}
\end{theorem}

$d(x,A)$ を見たら、任意の $a\in A$ に三角不等式を使ってから $a$ について下限を取る。
最小点の存在を問われたら、$a\mapsto d(x,a)$ の連続性と $A$ のコンパクト性を使う。

\begin{proof}
任意の $a\in A$ に三角不等式を使うと
$d(x,A)\le d(x,a)\le d(x,y)+d(y,a)$ である。
$a$ について下限を取り、$x,y$ を交換すれば (1) を得る。
(2) は、距離0なら各 $n$ で $d(x,a_n)<1/n$ となる $a_n\in A$ を選ぶことと、触点の定義から従う。
(3) は連続関数 $a\mapsto d(x,a)$ がコンパクト集合 $A$ 上で最小値を取ることによる。
\end{proof}

\begin{proposition}{コンパクト集合の一様な近傍}{tb-compact-neighborhood}
$K$ がコンパクトで $K\subset L$、$L$ が開なら、ある $r>0$ が存在して
\[
K_r=\{x\in X\mid d(x,K)<r\}\subset L
\]
となる。
\end{proposition}

各 $a\in K$ で $B(a;2r_a)\subset L$ を取り、$B(a;r_a)$ の有限部分被覆を選ぶ。
選ばれた半径の最小値を $r>0$ とすれば、$d(x,K)<r$ の点は
ある $a\in K$ の $B(a;2r_a)$ に入り、$x\in L$ となる。

\section{反例の最小セット}

連続写像について誤りやすい主張は、次の三例で処理できる。
\begin{enumerate}
  \item \textbf{開集合の像は開とは限らない。}
        $f:\R\to\R$, $f(x)=x^2$ で $f((-1,1))=[0,1)$。
  \item \textbf{コンパクト集合の逆像はコンパクトとは限らない。}
        定数写像 $f:\R\to\R$, $f(x)=0$ で $f^{-1}(\{0\})=\R$。
  \item \textbf{全射でも稠密集合の逆像は稠密とは限らない。}
        $f:\R\to\R$ を $\abs x\le1$ で0、$x>1$ で $x-1$、$x<-1$ で $x+1$ とする。
        これは連続全射だが、稠密集合 $\R\setminus\{0\}$ の逆像は
        $(-\infty,-1)\cup(1,\infty)$ であり稠密でない。
\end{enumerate}

「連続像が保つ」のはコンパクト性と連結性、
「連続写像の逆像が保つ」のは開集合と閉集合、と整理する。
